% ==============================================================================
% reports/frontmatter/abstract_vi.tex
% Tóm tắt báo cáo (Tiếng Việt)
% ==============================================================================

\chapter*{Tóm Tắt Báo Cáo}
\addcontentsline{toc}{chapter}{Tóm Tắt Báo Cáo}

Trong lĩnh vực an toàn thông tin, ứng dụng học máy trong các hệ thống phát hiện xâm nhập mạng (NIDS) thường đối mặt với hiện tượng đánh giá quá lạc quan khi phân tách dữ liệu ngẫu nhiên. Báo cáo bài tập lớn này thực hiện đánh giá thực nghiệm toàn diện bài toán phát hiện tấn công dò quét mật khẩu (Brute Force gồm \FTP\ và \SSH) trên tập dữ liệu CICIDS2017 theo trình tự thời gian, áp dụng kỹ thuật làm sạch ranh giới Purge ($60$ giây) và Embargo ($120$ giây) nhằm ngăn ngừa rò rỉ dữ liệu.

Nghiên cứu tập trung đánh giá khả năng tổng quát hóa từ biến thể đã học sang biến thể mới (kịch bản FTP $\to$ SSH): mô hình chỉ được huấn luyện trên tấn công \FTP\ buổi sáng và kiểm thử trên biến thể \SSH\ buổi chiều. Kết quả thực nghiệm trên ba thuật toán dạng cây (Cây quyết định, Rừng ngẫu nhiên, XGBoost) trong kịch bản giữ nguyên cổng dịch vụ cho thấy sự suy giảm hiệu năng nghiêm trọng: Rừng ngẫu nhiên chỉ đạt $\FOne = \RFTestFOne$ với tỷ lệ phát hiện SSH đạt $\RFTestSSHRecall$ (7/3,732 luồng), trong khi Cây quyết định và XGBoost hoàn toàn không phát hiện được tấn công SSH ($\FOne = 0.0000$). Phân tích khả năng giải thích bằng giá trị SHAP chỉ ra mô hình phụ thuộc áp đảo vào cổng dịch vụ đích (giá trị SHAP trung bình gấp \SHAPRatioTopOneTwo\ lần đặc trưng kế tiếp). Khi bóc tách đặc trưng cổng, tỷ lệ phát hiện SSH tăng nhẹ lên mức $0.19\% -- 0.38\%$ (F1 cao nhất đạt \RFTestWithoutPortFOne\ ở Rừng ngẫu nhiên), khẳng định cổng dịch vụ là lối tắt phân loại quan trọng nhưng sự khác biệt hình thái gói tin giữa hai giao thức cũng là yếu tố then chốt tạo nên sự dịch chuyển phân phối dữ liệu.

Trái lại, kịch bản phân tách ngẫu nhiên cho kết quả F1 rất cao ($0.9823 -- 0.9997$) do tương quan mạnh giữa các luồng trong cùng chiến dịch tấn công. Nghiên cứu cũng phân tích trường hợp điển hình về thiết kế pipeline: việc gộp tập kiểm định vào huấn luyện trước khi tìm kiếm siêu tham số đã vô tình đưa mẫu SSH vào mô hình, làm thay đổi bản chất bài toán và đẩy F1 lên \XGBRefitTestFOne. Từ đó, báo cáo đúc kết các bài học thực tiễn về thiết kế quy trình đánh giá khách quan và kiểm soát rò rỉ dữ liệu cho hệ thống phát hiện xâm nhập mạng.

\vspace{0.4cm}
\noindent\textbf{Từ khóa:} Hệ thống phát hiện xâm nhập (NIDS), CICIDS2017, Brute Force, Phân tách theo thời gian, Rò rỉ dữ liệu, Tổng quát hóa biến thể, SHAP.
